% ECONOMICS_PROBLEM: {"id": "C79-13", "title": "Determinacy of Dadaist Delta games", "jel_code": "C79", "source_id": "OP-0219"}
% !TeX program = xelatex
\documentclass[11pt,a4paper]{article}
\usepackage[margin=23mm]{geometry}
\usepackage{fontspec}
\setmainfont{Latin Modern Roman}
\usepackage{amsmath,amssymb,mathtools}
\usepackage{xcolor,enumitem,booktabs,longtable,array,xurl,hyperref}
\definecolor{muted}{HTML}{53636C}
\hypersetup{colorlinks=true,urlcolor=blue,linkcolor=blue}
\setlength{\parindent}{0pt}
\setlength{\parskip}{5pt}
\newcommand{\R}{\mathbb R}
\newcommand{\E}{\mathbb E}
\newcommand{\N}{\mathbb N}
\newcommand{\ind}{\mathbf 1}
\newcommand{\Var}{\operatorname{Var}}
\newcommand{\Cov}{\operatorname{Cov}}
\newcommand{\supp}{\operatorname{supp}}
\DeclareMathOperator*{\argmin}{arg\,min}
\DeclareMathOperator*{\argmax}{arg\,max}
\newcommand{\entrymeta}[3]{{\sffamily\small\color{muted}\textbf{\texttt{#1}}\quad|\quad #2\par #3\par}}
\newcommand{\Problemref}[1]{\texttt{#1}}
\newcommand{\JELref}[2]{\texttt{#2} (JEL \texttt{#1})}
\begin{document}
\section*{Determinacy of Dadaist Delta games}
\textbf{Problem ID:} \texttt{C79-13}\quad \textbf{JEL:} \texttt{C79}\par
% BEGIN ECONOMICS STATEMENT
\label{problem:OP-0219}
% GAME_THEORY_PROBLEM: {"local_id": "GT-089", "original_number": 89, "kind": "P", "entry_kind": "statement_only", "published": false, "review_required": true, "category": "game_theory"}
\label{game:GT-089}
{\sffamily\small\color{muted}
\noindent\textbf{ID:} GT-089. \textbf{Category:} Infinitary combinatorial games.
\textbf{Status:} P; Problem 3.4 in the inspected source.
\par}
\subsection*{Definitions and mathematical statement}
Work in ZFC. A combinatorial game form has designated Left and Right options and a well-founded option relation. Let $I$ be a set and let $(G_i)_{i\in I}$ be such forms, with their forms, rather than just their Conway-equivalence classes, fixed. The $\Delta$-sum is the indexed position
$\mathbf G=\sum_{i\in I}^{\Delta}G_i$. A Left (respectively Right) move replaces exactly one coordinate by a Left (respectively Right) option and leaves every other coordinate fixed.

Play may have ordinal length. At a limit stage the position is the coordinatewise eventual value of the earlier positions. This value exists because a fixed coordinate cannot undergo infinitely many moves along a well-founded option relation. After any limit stage the player designated First moves first again. More precisely, at the successor move numbered $\beta+j$, where $\beta=0$ or a limit ordinal and $j$ is a positive integer, First moves when $j$ is odd and Second when $j$ is even. A maximal legal alternating run ends when the player due to move has no legal move; that player loses. A strategy prescribes a legal move after each history at which its player must move; it is winning if every maximal compatible play is won by that player.

For a fixed choice $a\in\{\mathrm L,\mathrm R\}$ of First, let
$W_{\mathrm L}(\mathbf G,a)$ and $W_{\mathrm R}(\mathbf G,a)$ denote existence of a winning strategy for the indicated player. Ask whether
\[
 \forall I\ \forall (G_i)_{i\in I}\ \forall a\in\{\mathrm L,\mathrm R\},
 \qquad W_{\mathrm L}(\mathbf G,a)\ \lor\ W_{\mathrm R}(\mathbf G,a).
\]
Equivalently, does there exist a $\Delta$-sum and a choice of First for which both winning-strategy assertions fail? The countable-sum question restricts $I$ to $\mathbb N$.

\subsection*{Short English statement}
Under the specified transfinite play rule, must one player have a winning strategy for each choice of who starts?

\subsection*{Variants and scope boundaries}
The four usual outcomes $\mathrm N,\mathrm P,\mathrm L,\mathrm R$ are defined only when both starting-player choices are determined. Thus an undefined outcome need not mean failure of determinacy for both choices. In the default rule above, completed plays have winners; ``neither has a winning strategy'' is a determinacy failure, not a draw outcome. Stopping at time $\omega$ and declaring infinite play a draw, or changing the player who starts after a limit, defines a different game. The source also studies those alternative limit conventions separately.

\subsection*{Source relationship and status}
Definitions 3.1 and 3.3 fix the source's $\Delta$-sum and ordinary transfinite play. Problem 3.4 asks determinacy both for arbitrary set-indexed sums and for countable sums. This formulation makes explicit the starting-player quantifier and the source's assumption of Choice; finite normal-play determinacy alone does not settle it.

\subsection*{References}
\begingroup\footnotesize\raggedright
\noindent [S1] Paolo Lipparini, \emph{Infinite sums of combinatorial games (Dadaist games)} (2025), arXiv:2505.00614v1. Inspected locators: Section 3, Definitions 3.1 and 3.3, Problem 3.4; Sections 4--5 for other limit rules.
\url{https://arxiv.org/html/2505.00614v1}
\par\endgroup

\subsection*{Common conventions: Source mathematical conventions}

\textbf{Finite games.} For a finite set $A$, $\Delta(A)$ is its probability
simplex. Players choose independent mixed strategies in Nash equilibrium;
correlation is allowed only when explicitly part of the solution concept.
In a game with payoff functions $u_i$ and strategy profile $x$, an additive
$\varepsilon$ Nash equilibrium satisfies
\[
 u_i(x)\geq u_i(x_i',x_{-i})-\varepsilon
 \quad\text{for every player $i$ and every admissible deviation $x_i'$}.
\]
For numerical approximation, payoffs are normalized as locally specified.
An $\varepsilon$ payoff residual is distinct from distance at most
$\varepsilon$ to the set of exact equilibria. Point convergence, convergence
of empirical play, and convergence in distance to an equilibrium set are
also distinct.

\textbf{Computational representations.} Unless stated otherwise, a complexity
question takes rational data with binary encoding; polynomial time is measured
in total bit length. A fixed-accuracy polynomial algorithm is not necessarily
polynomial in $1/\varepsilon$, and neither is necessarily polynomial in
$\log(1/\varepsilon)$. Succinct games and value, demand, or sampling oracles
must declare their interfaces. Exact real solutions require an explicit
representation; an unspecified list of arbitrary real numbers is not a
finite computational output. A claimed algorithmic barrier is meaningful
only for its specified model and reductions.

\textbf{Dynamic games.} Histories, information, observation, and allowable
randomization are part of the model. A stationary strategy depends only on
the current state; a positional strategy in a deterministic graph game
chooses one action at each controlled vertex. Nash equilibrium at an initial
state and equilibrium after every history are different requirements.
An expectation of a pathwise $\liminf$ payoff, a $\liminf$ of expected averages,
a limiting expected average, and a uniform finite-horizon equilibrium are
different criteria. None is silently substituted for another.

\textbf{Allocation and mechanisms.} A complete allocation assigns every item
exactly once unless otherwise stated. Zero-valued goods, negative values,
capacity constraints, feasibility, lotteries, and copies of goods affect
fairness definitions and are specified locally. Dominant-strategy incentive
compatibility, Bayesian incentive compatibility, universal truthfulness,
and truthfulness in expectation impose different inequalities. Whenever
an objective ratio has zero denominator, its convention is stated or those
instances are excluded.

\textbf{Infinite-dimensional models.} Expectations, integrals, stochastic
equations, and optimization objectives require their local regularity,
integrability, filtrations, and admissible domains. A backward--forward
mean-field system is a boundary-value problem; it is not an initial-value
semiflow without an additional construction. Long-time, large-population,
and vanishing-noise limits require an order of limits and a convergence
topology. If a minimum need not be attained, the objective is its infimum
or supremum and the approximation requirement is stated separately.
% END ECONOMICS STATEMENT
\end{document}
